\documentclass[11pt]{article}

\usepackage{amsmath, amssymb, amsthm, geometry}
\usepackage{tikz-cd}
\usepackage{booktabs}
\usepackage{hyperref}
\usepackage{array}
\usepackage{longtable}
\usepackage{calc}
\usepackage{color}
\geometry{margin=1in}

\newtheorem{theorem}{Theorem}[section]
\newtheorem{proposition}[theorem]{Proposition}
\newtheorem{lemma}[theorem]{Lemma}
\newtheorem{definition}[theorem]{Definition}
\newtheorem{remark}[theorem]{Remark}
\newtheorem{conjecture}[theorem]{Conjecture}
\newtheorem{corollary}[theorem]{Corollary}
\newtheorem{axiom}[theorem]{Axiom}

\providecommand{\tightlist}{%
  \setlength{\itemsep}{0pt}\setlength{\parskip}{0pt}}

\begin{document}

\title{The Ultimate Axiom:\\
Deriving Quantum Mechanics from the Logic of Observation\\
{\large (Paper VI: The Necessity of $\mathbb{C}$)}}

\author{Zhou-Li Chen \\ co-nlang Research}
\date{May 2026}

\maketitle

\begin{abstract}
This paper completes the six-part series by demonstrating that the complex numbers $\mathbb{C}$---the foundation of quantum mechanics---emerge not as a postulate, but as a \textbf{necessary selection} within the framework of Sol\`{e}r's theorem. Building on the homological framework developed in Paper V, we prove that \textbf{within the class of division rings permitted by Sol\`{e}r's classification} ($\mathbb{R}$, $\mathbb{C}$, $\mathbb{H}$), any universe permitting non-trivial, continuous observation must necessarily select the complex field.

Our central result, \textbf{The Sol\`{e}r-Cohomology Theorem}, establishes this necessity in four rigorous steps:
\begin{enumerate}
\item Observation is contextual (no global section of the spectral presheaf).
\item Contextuality forces a non-split central extension, which in turn forces a non-trivial $d_2$ transgression in the Lyndon-Hochschild-Serre spectral sequence.
\item The transgression mechanism requires an abelian character group $\mathrm{Hom}(N, U(1)) \neq 0$.
\item The only compact, connected, abelian Lie group arising from a division ring in Sol\`{e}r's classification is $U(1)$, the phase group of $\mathbb{C}$.
\end{enumerate}

We show that $\mathbb{R}$ fails because its discrete phase group $O(1) \cong \mathbb{Z}/2$ cannot support continuous dynamics, and $\mathbb{H}$ fails because its non-abelian phase group $Sp(1) \cong SU(2)$ has trivial abelianization, annihilating the transgression mechanism. Only $\mathbb{C}$ survives.

Thus, quantum mechanics over the complex field is the unique mathematical structure within Sol\`{e}r's framework consistent with the possibility of meaningful observation.
\end{abstract}

%------------------------------------------------------------------
\section{Introduction: From Postulate to Theorem}
\label{introduction-from-postulate-to-theorem}
%------------------------------------------------------------------

\subsection{The Unreasonable Effectiveness of $\mathbb{C}$}
\label{the-unreasonable-effectiveness-of-c}

Quantum mechanics rests on the complex Hilbert space. The Schr\"{o}dinger equation, the uncertainty principle, entanglement---all require complex amplitudes. Yet for nearly a century, $\mathbb{C}$ has been treated as a \textbf{postulate}, justified only by empirical success.

\begin{quote}
Why complex numbers? Why not $\mathbb{R}$, $\mathbb{H}$, or some other division ring?
\end{quote}

Sol\`{e}r's theorem (1995) narrowed the possibilities to three: $\mathbb{R}$, $\mathbb{C}$, $\mathbb{H}$. But Sol\`{e}r's result is a \textbf{classification}, not a \textbf{selection principle}. It tells us what is possible, not which one permits non-trivial observation.

\subsection{The Thesis of Paper VI}
\label{the-thesis-of-paper-vi}

We establish that $\mathbb{C}$ is \textbf{selected} by a deeper principle: the requirement that ``observation'' be a meaningful operation. Specifically:

\textbf{Thesis}: Within the framework of Sol\`{e}r's classification, if a physical theory projects quantum systems onto their classical shadows (via the Bohrification functor), and if this observation is \textbf{contextual} (i.e., the spectral presheaf admits no global section), then the underlying field must necessarily be $\mathbb{C}$.

\subsection{The Argument in Brief}
\label{the-argument-in-brief}

Our derivation proceeds in four rigorous steps:
\begin{enumerate}
\item \textbf{Contextuality} $\rightarrow$ No global section of the spectral presheaf exists
\item \textbf{Cohomological Obstruction} $\rightarrow$ $H^2$ transgression is mandatory
\item \textbf{Phase Group Constraints} $\rightarrow$ $U(1)$ is the unique solution
\item \textbf{Division Ring Selection} $\rightarrow$ $\mathbb{C}$ is the only field supporting $U(1)$
\end{enumerate}

This paper unpacks each step, proving the Sol\`{e}r-Cohomology Theorem and thereby deriving quantum mechanics from first principles.

%------------------------------------------------------------------
\section{The Logic of Observation: From Contextuality to Cohomological Necessity}
\label{the-logic-of-observation}
%------------------------------------------------------------------

This section presents the rigorous mathematical derivation of quantum mechanics from the possibility of observation. We proceed in four steps, beginning with the physical requirement that observation be contextual.

%------------------------------------------------------------------
\subsection{Step 1: Contextuality as the Signature of Non-Trivial Observation}
\label{step-1-contextuality}

\textbf{The Setup}.
Recall from Paper~V~\cite{PaperV} the \textbf{Bohrification functor} $\mathcal{B}: \mathbf{Quant}^{\text{op}} \to \mathbf{Class}$, which projects quantum systems onto their classical shadows (the poset of maximal abelian subalgebras, MASAs). For a quantum algebra $A$, the spectral presheaf $\underline{\Sigma}$ over $\mathcal{B}(A)$ encodes the possible value assignments to observables in each classical context.

\textbf{The Physical Requirement}.
For observation to be physically meaningful in a quantum universe, it must capture the essential distinction between quantum and classical systems: the impossibility of assigning consistent values to all observables simultaneously. This is the content of the Kochen-Specker theorem.

\textbf{The Ultimate Axiom}.
We posit that observation is \textbf{contextual}:

\begin{axiom}[Contextuality of Observation]
\label{axiom:contextuality}
There exists at least one quantum algebra $A$ whose spectral presheaf $\underline{\Sigma}$ over the MASA poset $\mathcal{B}(A)$ admits \textbf{no global section}.
\end{axiom}

\begin{remark}[Equivalent Formulations]
By the Kochen-Specker theorem (formulated via Isham-Butterfield or Abramsky-Brandenburger~\cite{Abramsky2011}), the following are equivalent for a quantum algebra $A$:
\begin{enumerate}
\item The spectral presheaf $\underline{\Sigma}$ admits no global section
\item $A$ is non-commutative
\item The valuation $v: \text{Observables}(A) \to \mathbb{R}$ cannot be consistent across all contexts
\end{enumerate}
Thus, Axiom~\ref{axiom:contextuality} captures the essence of quantum contextuality.
\end{remark}

\begin{definition}[Quantum Contextuality]
\label{def:quantum-contextuality}
The absence of a global section for $\underline{\Sigma}$ is called \textbf{quantum contextuality}---the irreducible dependence of measurement outcomes on the choice of classical context.
\end{definition}

\textbf{Key Insight}.
Contextuality is not a bug but a feature: it is the signature that makes quantum mechanics non-classical. Where contextuality holds, classical logic fails to capture the quantum whole.

%------------------------------------------------------------------
\subsection{Step 2: From Contextuality to Cohomological Obstruction}
\label{step-2-cohomological-obstruction}

\textbf{The Group-Theoretic Translation}.
Consider a quantum system with symmetry group $G$ (e.g., the Pauli group) and center $N$ (the phase group). The MASA poset corresponds to the quotient $G/N$ (the abelianized classical phase space), while the full quantum structure corresponds to $G$.

\textbf{The Central Question}.
The existence of a global section for the spectral presheaf corresponds, via Heunen-Landsman-Spitters~\cite[Theorem 5.3]{Heunen2012}, to the algebraic question of whether the group extension:
\[
1 \longrightarrow N \longrightarrow G \longrightarrow G/N \longrightarrow 1
\]
\textbf{splits}---i.e., whether $G \cong N \rtimes G/N$ as a semidirect product.

\begin{theorem}[Contextuality Implies Non-Split Extension (Finite-Dimensional Group-Theoretic Models)]
\label{thm:contextuality-non-split}
Let $A = M_n(\mathbb{C})$ be a finite-dimensional quantum algebra whose symmetry group $G$ (e.g.\ the Pauli group $\mathcal{P}_n$) acts on $A$ by conjugation, with centre $N = Z(G)$. If the spectral presheaf $\underline{\Sigma}$ over the MASA poset $\mathcal{B}(A)$ admits no global section, then the associated central group extension
\[
1 \longrightarrow N \longrightarrow G \longrightarrow G/N \longrightarrow 1
\]
does \textbf{not} split. In particular, $G$ cannot be decomposed as a semidirect product $N \rtimes G/N$.
\end{theorem}

\begin{proof}
We proceed by contraposition in two stages: first for the algebraic structure of $A$, then for the group extension.

\textbf{Stage 1: From split extension to commutative algebra.}
Assume the group extension \textbf{splits}. Then there exists a group homomorphism (a section) $s: G/N \to G$ with $\pi \circ s = \mathrm{id}_{G/N}$. Because $N$ is central, the image $s(G/N)$ is a subgroup of $G$ isomorphic to $G/N$ and intersecting $N$ trivially. The group $G$ is then the internal direct product $G \cong N \times G/N$.

For the Pauli group and its finite-dimensional analogues, $G/N$ is abelian (indeed $G/N \cong (\mathbb{Z}/2)^{2n}$ for the $n$-qubit Pauli group). A split extension with abelian quotient therefore makes $G$ a central product of an abelian group with a phase group. The associated C*-algebra $A$ generated by this group representation is then commutative, because all generators commute up to a central phase that is absorbed into the direct product structure.

\textbf{Stage 2: From commutative algebra to global section.}
\cite[Theorem~5.3]{Heunen2012} proves: for a finite-dimensional C*-algebra $A$, the spectral presheaf $\underline{\Sigma}$ admits a global section iff $A$ is commutative. (The contrapositive: $M_n(\mathbb{C})$ with $n \ge 3$ has no global section.)

Combining the two stages:
\[
\text{split extension} \;\Rightarrow\; G \text{ central-direct-product} \;\Rightarrow\; A \text{ commutative} \;\Rightarrow\; \underline{\Sigma} \text{ has a global section}.
\]
By contraposition, if $\underline{\Sigma}$ admits no global section (Axiom~\ref{axiom:contextuality}), then the group extension cannot split.
\end{proof}

\begin{remark}[Scope of the Equivalence]
The proof above invokes the equivalence ``no global section $\Leftrightarrow$ non-commutative algebra'' from Heunen-Landsman-Spitters~\cite[Theorem 5.3]{Heunen2012}, which is rigorous for finite-dimensional C*-algebras. The bridge from ``non-commutative algebra'' to ``non-split group extension'' is explicit for group-theoretic models (Pauli group, Weyl-Heisenberg group) where the algebra is generated by a projective representation of $G/N$. For general infinite-dimensional C*-algebras, the relationship between topos-theoretic contextuality and group-cohomological splitting remains conjectural; see Paper V~\cite{PaperV} for the proposed functorial bridge.
\end{remark}

\begin{remark}[Finite vs. Infinite-Dimensional Scope]
The proof above invokes Heunen-Landsman-Spitters~\cite[Theorem 5.3]{Heunen2012}, which applies to finite-dimensional C*-algebras. The current result establishes the necessity of $\mathbb{C}$ for finite-dimensional quantum systems. The extension to infinite-dimensional systems requires additional analysis of whether the spectral presheaf admits global sections in the infinite-dimensional Bohrification framework; this is left for future work.
\end{remark}

\textbf{The LHS Spectral Sequence}.
To compute the obstruction to splitting explicitly, we now introduce the algebraic machinery that generates it from the group extension structure. The \textbf{Lyndon-Hochschild-Serre (LHS) spectral sequence} provides this framework. Specifically, the $d_2$ differential:
\[
d_2: E_2^{0,1} \longrightarrow E_2^{2,0}
\]
captures the transgression from the center's characters to the cohomology of the classical phase space.

\begin{theorem}[Non-Splitting Implies Non-Trivial Transgression]
\label{thm:non-splitting}
Consider a central extension $1 \to N \to G \to G/N \to 1$. The primary obstruction to splitting is the extension class $[e] \in H^2(G/N, N)$. Let $\chi: N \to U(1)$ be a continuous character of the centre. The induced map in cohomology,
\[
\chi_*: H^2(G/N, N) \longrightarrow H^2(G/N, U(1)),
\]
pushes the extension class forward to a derived obstruction class $[f_\chi] = \chi_*([e])$.

If the extension does not split, then $[e] \neq 0$. Moreover, there exists a character $\chi \in \mathrm{Hom}(N, U(1))$ such that its image under the LHS transgression differential is non-trivial:
\[
d_2(\chi) = [f_\chi] \neq 0 \in H^2(G/N, U(1)).
\]
\end{theorem}

\begin{proof}
For a central extension, the primary splitting obstruction is the extension class $[e] \in H^2(G/N, N)$. Standard group cohomology identifies: the extension splits if and only if $[e] = 0$.

Now consider the LHS spectral sequence with coefficients in $U(1)$ (viewed as a trivial $G$-module). The $E_2$-page contains the term $E_2^{0,1} = H^0(G/N, H^1(N, U(1)))$. Because $N$ is central, the conjugation action of $G/N$ on $N$ is trivial; hence $G/N$ acts trivially on $\mathrm{Hom}(N, U(1))$, and $E_2^{0,1} \cong \mathrm{Hom}(N, U(1))$.

The $d_2$ transgression map is:
\[
d_2: E_2^{0,1} = \mathrm{Hom}(N, U(1)) \longrightarrow E_2^{2,0} = H^2(G/N, U(1)).
\]
A standard diagram chase (see Brown~\cite[Chapter VII, \S 6]{Brown82}) shows that for any character $\chi: N \to U(1)$, the transgression $d_2(\chi)$ equals the pushforward of the extension class:
\[
d_2(\chi) = \chi_*([e]) \in H^2(G/N, U(1)).
\]
If the extension does not split, then $[e] \neq 0$. Because $N = U(1)$ (the case we shall establish for $\mathbb{C}$) is a compact abelian Lie group, its character group $\mathrm{Hom}(U(1), U(1)) \cong \mathbb{Z}$ is rich enough to detect non-trivial classes: the identity character $\chi_0(z) = z$ yields $d_2(\chi_0) = \chi_{0,*}([e]) = [e]$ (under the identification $N = U(1)$ with coefficients in $U(1)$). Hence $d_2(\chi_0) = [e] \neq 0$.

More generally, even before identifying $N$ with $U(1)$, if $N$ is any compact abelian group with $[e] \neq 0$, Pontryagin duality guarantees the existence of a character $\chi$ with $\chi_*([e]) \neq 0$. Thus non-splitting always produces a non-trivial derived obstruction $[f_\chi] = d_2(\chi)$ for some $\chi$.
\end{proof}

\begin{corollary}[Kochen-Specker as Necessity]
\label{cor:ks-necessity}
The contextuality of observation (no global section for the spectral presheaf) necessarily implies the existence of a non-vanishing obstruction class $[f] \in H^2(G/N, U(1))$, realized as the pushforward $[f] = \chi_*([e])$ of the primary extension class $[e]$ via a character $\chi$ of the centre. This class is precisely the Kochen-Specker contextuality of Paper III~\cite{PaperIII}.
\end{corollary}

\textbf{The Chain So Far}.
We have established:
\[
\boxed{\text{No global section} \Rightarrow \text{Non-split extension} \Rightarrow d_2 \neq 0 \Rightarrow H^2(G/N, U(1)) \neq 0}
\]

The nature of the phase group remains to be determined. The Kochen-Specker paradox is not an accident---it is forced by the logic of observation.

%------------------------------------------------------------------
\subsection{Step 3: The Pincer Movement --- Locking $U(1)$}
\label{step-3-pincer-movement}

We have proven that non-trivial observation implies a non-trivial obstruction in $H^2(G/N, \text{Center})$. But what is this ``Center''? The nature of the coefficient group remains undetermined.

\textbf{The Question}.
The cohomology class $[f] = d_2(\chi)$ requires a coefficient group $Z(G)$ (the center of the quantum symmetry). We know $Z(G)$ must be the phase group $N$. But which groups can serve as $N$ for a division ring in Sol\`{e}r's classification?

Candidates:
\begin{itemize}
\item \textbf{$\mathbb{R}$}: Phase group is $O(1) = \{\pm 1\} \cong \mathbb{Z}/2$ (discrete)
\item \textbf{$\mathbb{C}$}: Phase group is $U(1)$ (continuous, abelian)
\item \textbf{$\mathbb{H}$}: Phase group is $Sp(1) \cong SU(2)$ (continuous, non-abelian)
\end{itemize}

Only one of these supports the continuous, abelian structure required by geometric phases. We now eliminate the alternatives through a \textbf{pincer movement}: physical geometry closes one door, categorical homological algebra closes the other.

%------------------------------------------------------------------
\subsubsection{The Physical-Geometric Boundary: Excluding $\mathbb{R}$}
\label{excluding-r}

\textbf{Algebraic Premise}.
If the underlying division ring were $\mathbb{R}$, its unitary group (phase group) degenerates to the orthogonal group $O(1) = \{\pm 1\} \cong \mathbb{Z}/2$.

\textbf{Physical and Geometric Contradiction}.
In Papers I and II~\cite{PaperI,PaperII}, we proved that the geometric shadow of quantum observation manifests as flat line bundles over classical phase space $M$, classified by $\check{H}^1(M, \text{Phase Group})$.

If the phase group were $\mathbb{Z}/2$, then $\check{H}^1(M, \mathbb{Z}/2)$ could only describe discrete topological defects (such as orientation flips \`{a} la M\"{o}bius strips). However, physical phenomena---including the Aharonov-Bohm effect and Berry phase---require phases to \textbf{accumulate continuously} along paths.

More precisely: a continuous one-parameter family of dynamics, as generated by a Hamiltonian $H$ via $U(t) = e^{-iHt}$, requires a continuous homomorphism $\mathbb{R} \to P$ from the real line (time) to the phase group $P$. Since $\mathbb{R}$ is connected, the image of this homomorphism must lie in the connected component of the identity in $P$. If $P = \{\pm 1\}$ is discrete, the only connected subgroup is $\{1\}$; every continuous homomorphism $\mathbb{R} \to \{\pm 1\}$ is trivial. Consequently, no non-trivial continuous time evolution is possible. This makes Stone's theorem inapplicable: there is no faithful continuous unitary representation of time translation.

Equivalently, from a topological perspective:
\begin{itemize}
\item The fundamental group of $U(1)$ is $\mathbb{Z}$, allowing winding numbers and continuous holonomy
\item The group $\{\pm 1\}$ has trivial fundamental group, forbidding continuous phase accumulation
\end{itemize}

\begin{theorem}[Exclusion of $\mathbb{R}$ --- The Continuity Constraint]\label{thm:exclude-r}
If the phase group is discrete ($\mathbb{Z}/2$), the theory cannot support continuous geometric phases or unitary time evolution. Therefore, $\mathbb{R}$ is insufficient for quantum mechanics.
\end{theorem}

\textbf{Conclusion}.
To allow the $d_2$ obstruction (contextuality) to coexist with continuous evolution, the phase group must be a \textbf{continuous Lie group}.

%------------------------------------------------------------------
\subsubsection{The Categorical-Homological Boundary: Excluding $\mathbb{H}$}
\label{excluding-h}

This step is the soul of the argument, directly echoing the adjunction framework.

\textbf{Algebraic Premise}.
If the underlying division ring were $\mathbb{H}$, its unitary group extends to $Sp(1) \cong SU(2)$, a \textbf{non-abelian} continuous Lie group.

\textbf{Categorical Contradiction: Collapse of MASAs}.
The Bohrification functor $\mathcal{B}$ projects quantum systems onto the poset of \textbf{maximal abelian subalgebras (MASAs)}.

If phase patching were driven by $SU(2)$, transitions between local classical contexts would themselves be non-commutative. This means the ``phase'' concept cannot reside stably in the algebra's center. If phases do not commute, they cannot commute with all observables in a MASA---\textbf{directly violating the ``maximal abelian'' definition}.

\textbf{Why the Coefficient Module Must Be $U(1)$}.
Before invoking the LHS spectral sequence, we must justify the choice of coefficient module. The obstruction class $[f]$ arises from the attempt to lift classical data (abelian characters on $G/N$) to quantum data (a representation of the full group $G$). In quantum mechanics, the ``quantum data'' are probability amplitudes, and the only physically observable quantity derived from an amplitude is its modulus squared (the Born rule). Therefore, two amplitudes that differ by a phase factor of unit modulus yield identical predictions. The set of such phase factors, for \emph{any} normed division algebra, is precisely the group of unit-norm elements. For $\mathbb{C}$, this is $U(1)$; for $\mathbb{H}$, it is $Sp(1) \cong SU(2)$.

However, the \emph{physical equivalence} of amplitudes under phase is an \textbf{abelian} relation: if $\psi \sim e^{i\theta}\psi$ and $\phi \sim e^{i\phi}\phi$, the combined system satisfies $\psi \otimes \phi \sim e^{i(\theta+\phi)}(\psi \otimes \phi)$. The phase redundancy is always additive, regardless of the underlying division ring. Consequently, the \emph{physically relevant} character group that measures this redundancy is the group of homomorphisms from the phase group $N$ to the circle group $U(1)$, because $U(1)$ is the universal recipient of abelian phase information. It is in this sense---not by presupposing $\mathbb{C}$---that $U(1)$ appears as the coefficient module.

\textbf{Homological Contradiction: Collapse of LHS}.
With the coefficient module fixed as $U(1)$, the LHS spectral sequence has $E_2^{0,1}$ term $H^0(G/N, H^1(N, U(1))) \cong \mathrm{Hom}(N, U(1))^{G/N}$. For the $d_2$ transgression to have a non-trivial domain, we require $\mathrm{Hom}(N, U(1)) \neq 0$.

If $N = SU(2)$, this group is \textbf{perfect}: $[SU(2), SU(2)] = SU(2)$.\linebreak Consequently $\mathrm{Hom}(SU(2), U(1)) = 0$ (any homomorphism to abelian $U(1)$ factors through $SU(2)^{\text{ab}} = \{1\}$). This annihilates the $d_2$ mechanism regardless of whether $N$ is central.

\begin{theorem}[Exclusion of $\mathbb{H}$ --- The Abelian Constraint]\label{thm:exclude-h}
If the underlying division ring is $\mathbb{H}$, its unitary phase group is $Sp(1) \cong SU(2)$. Because $SU(2)$ admits no non-trivial continuous characters into $U(1)$, the Lyndon-Hochschild-Serre (LHS) $d_2$ transgression mechanism---which requires a non-trivial character group $\mathrm{Hom}(N, U(1))$---collapses. Kochen-Specker contextuality, which is realized cohomologically as a non-trivial $d_2$ transgression, is therefore impossible under a quaternionic phase group. Hence $\mathbb{H}$ is incompatible with the observation framework.
\end{theorem}

\begin{proof}
Let $N$ be the phase group of the division ring under consideration. The LHS spectral sequence for the group extension $1 \to N \to G \to G/N \to 1$ evaluates the cohomology of $G$ with coefficients in $U(1)$. As argued above, $U(1)$ is the physically motivated coefficient module because it universally encodes abelian phase redundancy.

\begin{enumerate}
\item \textbf{The $E_2$ Page Setup}: The mechanism that generates Kochen-Specker contextuality is the $d_2$ transgression differential:
\[
d_2: E_2^{0,1} \longrightarrow E_2^{2,0}
\]
The term $E_2^{0,1}$ is $H^0(G/N, H^1(N, U(1)))$.

\item \textbf{Evaluating $H^1(N, U(1))$}: By the definition of first group cohomology with trivial coefficients, $H^1(N, U(1)) \cong \mathrm{Hom}(N, U(1))$, the group of continuous characters from $N$ to $U(1)$. Since $N$ is central, $G/N$ acts trivially on $U(1)$, so $E_2^{0,1} \cong \mathrm{Hom}(N, U(1))$.

\item \textbf{The Abelianization of $SU(2)$}: For $N = SU(2)$, any character $\chi: SU(2) \to U(1)$ must factor through the abelianization $SU(2)^{\text{ab}} = SU(2)/[SU(2), SU(2)]$.

\item \textbf{The Topological Collapse}: The Lie group $SU(2)$ is semisimple and simply-connected. Consequently, $[SU(2), SU(2)] = SU(2)$ (algebraically, $SU(2)$ is a \textit{perfect group}). Thus $SU(2)^{\text{ab}} \cong \{1\}$.

\item \textbf{Vanishing of Characters}: Because the abelianization is trivial:
\[
\mathrm{Hom}(SU(2), U(1)) = 0
\]
This strictly forces $E_2^{0,1} = 0$.

\item \textbf{Conclusion}: Since the domain of the $d_2$ differential is the trivial group ($0$), the $d_2$ transgression map is identically zero. Without a non-trivial $d_2$ map, it is homologically impossible to generate the obstruction class $[f] \in H^2(G/N, U(1))$. The Kochen-Specker paradox---the signature of non-trivial observation derived in Theorem~\ref{thm:non-splitting}---\textbf{cannot exist} if the phase group is $SU(2)$.
\end{enumerate}

Since Axiom~\ref{axiom:contextuality} demands contextuality of observation (and thus a non-vanishing $d_2$ obstruction), the phase group $N$ cannot be $SU(2)$. Therefore, the underlying division ring cannot be $\mathbb{H}$.
\end{proof}

\textbf{Conclusion}.
The non-abelian phase group of $\mathbb{H}$ destroys the abelian foundation of the $\mathbf{Class}$ category in the conjectured $\mathcal{Q} \dashv \mathcal{B}$ adjunction (see Paper V~\cite{PaperV}). Therefore, the phase group must be \textbf{abelian}.

%------------------------------------------------------------------
\subsubsection{The Unique Solution Emerges: $\mathbb{C}$ and $U(1)$}
\label{unique-solution}

Synthesizing the two boundaries, a universe maintaining ``continuous observation'' and ``local abelian logic'' requires a phase group satisfying:
\begin{enumerate}
\item \textbf{Continuous} (excluding $\mathbb{R}$)
\item \textbf{Abelian} (excluding $\mathbb{H}$)
\item \textbf{Compact} (for bounded probabilities)
\end{enumerate}

\begin{corollary}[Phase Group Constraints]\label{cor:phase-constraints}
The phase group must be a compact, connected, abelian Lie group.
\end{corollary}

\begin{proof}
\textbf{Compactness} is required for bounded probabilities: the phase group acts on normalized states, and its orbits must be compact to ensure finite probability measures.

\textbf{Abelian} is required by Theorem~\ref{thm:exclude-h}: a non-abelian phase group annihilates the LHS $d_2$ transgression mechanism, rendering contextuality impossible.

\textbf{Connected} is required by continuous dynamics. Stone's theorem guarantees that any continuous one-parameter unitary group $U(t)$ has the form $U(t) = e^{-iHt}$ for a self-adjoint Hamiltonian $H$. The map $t \mapsto U(t)$ is a continuous path in the phase group starting at the identity. If the phase group were disconnected, the image of this path would be confined to the connected component of the identity, but transitions between different connected components would require discontinuous jumps, violating the smoothness of Hamiltonian evolution. More precisely, any one-parameter subgroup $\mathbb{R} \to P$ of the phase group $P$ must factor through the connected component $P_0$, because $\mathbb{R}$ is connected. If $P$ were disconnected (e.g. $O(2)$), the continuous dynamics would never access the other components, rendering them physically irrelevant; the effective phase group would be $P_0$ anyway.
\end{proof}

\begin{theorem}[Classification]\label{thm:classification}
The only compact, connected, abelian Lie group is a torus $T^n = (S^1)^n$.
\end{theorem}

\begin{proof}
By the classical classification of compact connected abelian Lie groups~\cite{Hochschild1965}.
\end{proof}

\begin{corollary}[Why $n=1$ for Division Rings]\label{cor:n-equals-1}
For a division ring's unitary group within Sol\`{e}r's classification, we require $n=1$, yielding $U(1) = S^1$.
\end{corollary}

\begin{proof}
By the Frobenius theorem on real division algebras~\cite{Baez2002}, the only finite-dimensional associative division algebras over $\mathbb{R}$ are $\mathbb{R}$, $\mathbb{C}$, and $\mathbb{H}$. Their respective centers and unitary groups are:

\begin{itemize}
\item $\mathbb{R}$: Center is $\mathbb{R}$ itself, unitary group is $O(1) \cong \{\pm 1\}$ (discrete, excluded by Theorem~\ref{thm:exclude-r})
\item $\mathbb{H}$: Center is $\mathbb{R}$, unitary group (as quaternions of unit norm) is $Sp(1) \cong SU(2)$ (non-abelian, excluded by Theorem~\ref{thm:exclude-h})
\item $\mathbb{C}$: Center is $\mathbb{C}$ itself, unitary group is $U(1) \cong S^1$ ($n=1$)
\end{itemize}

Among Sol\`{e}r's three candidates, only $\mathbb{C}$ yields a compact, connected, abelian phase group, and it is necessarily the one-dimensional torus $U(1) = T^1$.
\end{proof}

And $U(1)$ can only exist as the phase group of the complex field $\mathbb{C}$.

%------------------------------------------------------------------
\subsection{Step 4: The Uniqueness of $\mathbb{C}$}
\label{step-4-uniqueness}

\textbf{The Division Ring Analysis}:

\begin{center}
\begin{tabular}{ccc}
\toprule
\textbf{Division Ring} & \textbf{Phase Group} & \textbf{Verdict} \\
\midrule
$\mathbb{R}$ & $O(1) = \{\pm 1\}$ (discrete) & \textcolor{red}{\texttimes} Excluded by Thm~\ref{thm:exclude-r} \\
$\mathbb{H}$ & $Sp(1) \cong SU(2)$ (non-abelian) & \textcolor{red}{\texttimes} Excluded by Thm~\ref{thm:exclude-h} \\
$\mathbb{C}$ & $U(1)$ (continuous, abelian) & \textcolor{green}{\checkmark} \textbf{The unique solution} \\
\bottomrule
\end{tabular}
\end{center}

\textbf{The Sol\`{e}r-Cohomology Theorem}.
\[
\text{Observation} \rightarrow \text{Contextuality} \rightarrow H^2\text{-Transgression} \rightarrow U(1) \rightarrow \mathbb{C}
\]

Each implication is necessary:
\begin{enumerate}
\item Contextuality (no global section) $\Rightarrow$ non-split extension (Theorem~\ref{thm:contextuality-non-split})
\item Non-split extension $\Rightarrow$ non-trivial $d_2$ transgression (Theorem~\ref{thm:non-splitting})
\item Non-trivial transgression $\Rightarrow$ requires $U(1)$ phase group (Theorems~\ref{thm:exclude-r} and~\ref{thm:exclude-h})
\item $U(1)$ $\Rightarrow$ requires $\mathbb{C}$ as division ring (Sol\`{e}r classification)
\end{enumerate}

\begin{theorem}[The Sol\`{e}r-Cohomology Theorem]\label{thm:soler-cohomology}
\textbf{Part I (Within Sol\`{e}r's Classification).} Let a physical theory be founded on an orthomodular space over a division ring $K \in \{\mathbb{R}, \mathbb{C}, \mathbb{H}\}$ (as permitted by Sol\`{e}r's theorem). If the theory supports non-trivial observation (contextuality implies non-split extension implies $d_2 \neq 0$) and continuous dynamic evolution, then $K$ must be isomorphic to $\mathbb{C}$.

Equivalently: \textbf{Within Sol\`{e}r's classification, $\mathbb{C}$ is the unique division ring permitting cohomologically non-trivial observation.}
\end{theorem}

\begin{remark}[Associativity as an implicit boundary]
Sol\`{e}r's theorem itself restricts to associative division algebras. The octonions $\mathbb{O}$ fall outside this classification precisely because they are non-associative. Our entire framework---the LHS spectral sequence, group cohomology, and the category-theoretic observation adjunction---relies on associativity at every step. Thus associativity is not an additional assumption but a structural prerequisite embedded in Sol\`{e}r's hypothesis. The exclusion of $\mathbb{O}$ is therefore automatic: the theorem simply does not apply to non-associative algebras. See Section~\ref{octonionic-h3} for a discussion of whether higher cohomology ($H^3$) might accommodate non-associative structures.
\end{remark}

\begin{corollary}[Necessity of Quantum Mechanics within Sol\`{e}r's Framework]\label{cor:necessity}
Quantum mechanics over $\mathbb{C}$ is the unique structure within Sol\`{e}r's classification consistent with non-trivial, continuous observation. The complex numbers are not merely one option among three---they are \textbf{selected} by the requirements of observation.
\end{corollary}

\textbf{The Complete Chain}:
\[
\boxed{\text{Observation} \Rightarrow \text{Contextuality} \Rightarrow H^2\text{-Transgression} \Rightarrow U(1) \Rightarrow \mathbb{C}}
\]

%------------------------------------------------------------------
\begin{remark}[The Sol\`{e}r-EML Correspondence]\label{rem:soler-eml}
The derivation of $\mathbb{C}$ from the logic of observation has a suggestive computational echo. In the Epilogue~\cite{Epilogue}, we introduced the EML system $\text{eml}(x, y) = \exp(x) - \ln(y)$. Over $\mathbb{C}$, the evaluation $\ln(-1) = i\pi$ yields a well-defined $U(1)$ phase ambiguity. Over $\mathbb{H}$, the logarithm of a negative scalar admits a continuous $S^2$ of pure imaginary unit solutions, producing a non-abelian ``spherical branch ambiguity'' that prevents consistent single-valued computation. This computational obstruction mirrors the homological obstruction $[f] \in H^2(G/N, U(1))$: both require an abelian phase group. We leave a rigorous formulation of this correspondence---and whether computational completeness over a division algebra is equivalent to the existence of a non-trivial LHS transgression---as a question for future work (see Section~\ref{eml-open-question}).
\end{remark}

%------------------------------------------------------------------
\section{Philosophical Implications}
\label{philosophical-implications}
%------------------------------------------------------------------

\subsection{The End of the Postulate}
\label{the-end-of-the-postulate}

Quantum mechanics has traditionally been founded on postulates:
\begin{enumerate}
\item States are rays in a complex Hilbert space
\item Observables are Hermitian operators
\item Evolution is unitary
\item Measurement projects onto eigenstates
\end{enumerate}

The Sol\`{e}r-Cohomology Theorem eliminates Postulate 1. Within Sol\`{e}r's classification, the complex Hilbert space is not arbitrary---it is \textbf{selected} by the requirements of non-trivial, continuous observation.

\subsection{Logic as Geometry}
\label{logic-as-geometry}

Our result vindicates the Leibnizian dream: the laws of physics are not contingent, but necessary. The complex numbers, far from being a mathematical convenience, are forced by the logic of observation itself.

In this view:
\begin{itemize}
\item \textbf{Logic} (the structure of observation)
\item \textbf{Geometry} (the topology of the MASA nerve)
\item \textbf{Algebra} (the field of amplitudes)
\end{itemize}

are three faces of a single necessity.

\subsection{Counterfactual Physics}
\label{counterfactual-physics}

What would a universe with $\mathbb{R}$ or $\mathbb{H}$ look like?

\textbf{The $\mathbb{R}$-universe}: Discrete, classical. No interference, no entanglement, no quantum computation. A deterministic clockwork.

\textbf{The $\mathbb{H}$-universe}: Non-abelian phase structure. No stable MASAs, no consistent classical contexts. The phase group $SU(2)$ forces transitions between contexts to be non-commutative, preventing the existence of a stable abelian center required for consistent classical observation.

Only the $\mathbb{C}$-universe permits the delicate balance: enough structure for quantum phenomena, enough regularity for classical emergence.

%------------------------------------------------------------------
\section{Connections to Earlier Papers}
\label{connections-to-earlier-papers}
%------------------------------------------------------------------

\subsection{Paper I--II: Why Geometric Phases Require $\mathbb{C}$}
\label{paper-i-ii}

The Berry phase and Chern classes require continuous holonomy. This is only possible with $U(1)$ coefficients, which only exist over $\mathbb{C}$.

\subsection{Paper III: Why KS Contextuality Requires $\mathbb{C}$}
\label{paper-iii}

The Peres-Mermin square's sign obstruction is a $d_2$ transgression. For this to yield consistent probability violations (not mere discrete paradoxes), the coefficient group must be $U(1)$, forcing $\mathbb{C}$.

\subsection{Paper IV--V: The Hierarchy of Division Rings}
\label{paper-iv-v}

Our framework explains the division ring hierarchy:
\begin{itemize}
\item $\mathbb{R}$: Too rigid (discrete phase group)
\item $\mathbb{H}$: Too chaotic (non-abelian phase group)
\item $\mathbb{C}$: The Goldilocks zone
\end{itemize}

\subsection{Paper VI: The Completion}
\label{paper-vi-completion}

Paper VI completes the circle: from the possibility of observation to the necessity of $\mathbb{C}$, from the logic of contextuality to the geometry of the quantum world.

%------------------------------------------------------------------
\section{Open Questions (For Future Work)}
\label{open-questions}
%------------------------------------------------------------------

\subsection{Can We Eliminate Other Postulates?}
\label{eliminate-other-postulates}

If Postulate 1 (complex Hilbert space) can be derived, can Postulates 2--4 also be derived from the logic of observation? This is the program of \textbf{operational quantum mechanics}.

\subsection{The Nature of Contextuality}
\label{nature-of-contextuality}

Our derivation assumes the contextuality of observation (no global section for the spectral presheaf). Can we prove this from even more primitive axioms? Or is contextuality itself the ultimate primitive?

\subsection{Beyond the Associative Boundary: The Octonionic $H^3$ Frontier}
\label{octonionic-h3}

\textit{Having established $\mathbb{C}$ as the unique stable domain of observation, we ask: does the universe permit breaking the associative boundary under extreme conditions?}

\textbf{The Topological Shadow of Non-Associative Observation}.
If, in frameworks beyond the Standard Model (such as M-theory), observation were to permit transient non-associativity, this anomaly would not manifest in $H^2$ (non-commutativity) but necessarily in \textbf{$H^3$}.

\textbf{Rendezvous with Paper IV}.
This is precisely the \textbf{Borromean Contextuality} foreshadowed in Paper IV~\cite{PaperIV}. When three pairwise compatible contexts fail to compose associatively at the global level, a non-trivial Dixmier-Douady class emerges.

\begin{conjecture}[The Exceptional Horizon Conjecture]\label{conj:exceptional-horizon}
The Kochen-Specker paradox ($H^2$) of standard quantum mechanics is a phenomenon over the complex field $\mathbb{C}$; higher-order Borromean anomalies ($H^3$) are topological shadows projected by the octonions $\mathbb{O}$ (or the exceptional Jordan algebra $\mathfrak{h}_3(\mathbb{O})$) into the physical world. This suggests that at the Planck scale or in the depths of string theory, the logic of observation may require a non-associative geometry.
\end{conjecture}

\subsection{Beyond Sol\`{e}r: Deriving Orthomodularity}
\label{beyond-soler}

Sol\`{e}r's theorem assumes an orthomodular lattice structure. Can we derive orthomodularity itself from the conjectured observation adjunction $\mathcal{Q} \dashv \mathcal{B}$ (Paper V~\cite{PaperV})? This would be the \textbf{ultimate} derivation---quantum logic from pure category theory.

\subsection{The EML-Quantum Parallel}
\label{eml-open-question}

In Remark~\ref{rem:soler-eml} we noted that the EML system's ability to evaluate $\ln(-1)$ over $\mathbb{C}$ mirrors the quantum $U(1)$ phase ambiguity. A rigorous proof that ``EML completeness implies $\mathbb{C}$'' would require developing a non-abelian logarithmic multi-valuedness theory for $\mathbb{H}$ and analyzing the corresponding obstruction structure. Is computational completeness over a division algebra exactly coextensive with the existence of a non-trivial LHS transgression?

%------------------------------------------------------------------
\section{Conclusion: The Necessity of It All}
\label{conclusion}
%------------------------------------------------------------------

We began with a question: Why complex numbers?

We end with an answer: Because observation must be possible.

Within the framework of Sol\`{e}r's theorem, the complex numbers are not a contingent feature of our universe, but the unique selection consistent with non-trivial, continuous observation. If there is a universe in which quantum systems can be observed within Sol\`{e}r's classification, that universe must run on complex amplitudes.

Quantum mechanics is not merely true. Within its logical framework, it is \textbf{necessarily} true.

%------------------------------------------------------------------
\subsection{The Bootstrap of Observation: An Escherian Closing}
\label{bootstrap-of-observation}

\textit{Traditional derivations seek an absolute starting point---a first cause from which all else follows. But when we reach the deepest layer of cosmic logic, we find not a linear chain of deduction, but a perfect self-consistent loop.}

\textbf{The Chicken and the Egg, Resolved}.
One might ask: Does the logic of observation \textit{create} the complex numbers? Or do the complex numbers \textit{enable} observation? This paper has walked a careful line, proving that within Sol\`{e}r's classification, only $\mathbb{C}$ permits cohomologically non-trivial observation. Yet one cannot deny that our very definition of ``quantum algebra'' in Section~\ref{step-1-contextuality} already presupposes the structure of $M_n(\mathbb{C})$.

The Sol\`{e}r-Cohomology Theorem reveals a deeper truth: \textbf{``meaningful observation'' and ``complex cohomology'' are co-emergent}. They do not stand in causal sequence, but in adjunction---like the conjectured functors $\mathcal{Q} \dashv \mathcal{B}$ themselves (Paper V~\cite{PaperV}). There is no ``first.'' There is only the elegant mutual necessity of complementary structures.

\textbf{The Universe as Bootstrap System}.
The physicist's dream has always been to find the ultimate axiom from which all laws descend. Paper VI suggests a different picture: the universe as a \textbf{bootstrap system}, a M\"{o}bius strip of logic where:
\begin{itemize}
\item The possibility of observation requires $H^2$-transgression
\item $H^2$-transgression requires $U(1)$
\item $U(1)$ requires $\mathbb{C}$
\item And $\mathbb{C}$ is precisely the structure that makes meaningful observation possible
\end{itemize}

This is not circular reasoning in the pejorative sense. It is \textbf{circular coherence}---the mark of a system that is complete unto itself, needing no external scaffolding.

\textbf{Escher's Hands, Drawing Each Other}.
Like the lithograph of two hands sketching one another into existence, the logic of observation and the geometry of the complex field co-create the reality we inhabit. Sol\`{e}r's theorem gave us three candidates; the double boundary (continuous, abelian) selected one. But this selection was not imposed from outside---it was the only self-consistent knot in the fabric of mathematical physics.

\begin{quote}
\textit{``What we cannot speak about we must pass over in silence.''} --- Wittgenstein

\textit{``What we can observe, we must describe in complex amplitudes.''}\linebreak --- The Sol\`{e}r-Cohomology Theorem
\end{quote}

The six-paper journey ends here, not with a foundation, but with a \textbf{mirroring}---the deepest kind of certainty mathematics can offer.

%------------------------------------------------------------------
\begin{thebibliography}{9}

\bibitem[Soler1995]{Soler1995}
M.~P. Sol\`{e}r,
``Characterization of Hilbert spaces with orthomodular spaces,''
\textit{Communications in Algebra} \textbf{23} (1995), 219--243.

\bibitem[PaperI]{PaperI}
Chen, Z.-L.
From Contextuality to Phase Cohomology:
A Computable Bridge Between Bohrification and Geometric Quantization.
\textit{Zenodo, DOI: 10.5281/zenodo.20072818}, 2026.

\bibitem[PaperII]{PaperII}
Chen, Z.-L.
Semiclassical Reconstruction of Riemann Surfaces from Bohrification.
\textit{Zenodo, DOI: 10.5281/zenodo.20073010}, 2026.

\bibitem[PaperIII]{PaperIII}
Chen, Z.-L.
Kochen--Specker Contextuality as Central Extension: The Peres--Mermin Square and the Pauli Group.
\textit{Zenodo, DOI: 10.5281/zenodo.20073127}, 2026.

\bibitem[PaperIV]{PaperIV}
Chen, Z.-L.
The Cohomological Obstruction Ladder: Lyndon--Hochschild--Serre Transgression and the $H^3$ Frontier.
\textit{Zenodo, DOI: 10.5281/zenodo.20073184}, 2026.

\bibitem[PaperV]{PaperV}
Chen, Z.-L.
Observation as Functor: 
The Adjunction of Quantum and Classical.
\textit{Zenodo, DOI: 10.5281/zenodo.20073318}, 2026.

\bibitem[Epilogue]{Epilogue}
Chen, Z.-L.
The Algebraic Logic of Geometry.
\textit{Zenodo, DOI: 10.5281/zenodo.20073253}, 2026.

\bibitem[Heunen2012]{Heunen2012}
C.~Heunen, N.~P. Landsman, and B.~Spitters,
``Bohrification of operator algebras and quantum logic,''
\textit{Synthese} \textbf{186} (2012), 719--752.

\bibitem[Abramsky2011]{Abramsky2011}
S.~Abramsky and A.~Brandenburger,
``The sheaf-theoretic structure of non-locality and contextuality,''
\textit{New Journal of Physics} \textbf{13} (2011), 113036.

\bibitem[Baez2002]{Baez2002}
J.~C. Baez,
``The octonions,''
\textit{Bulletin of the American Mathematical Society} \textbf{39} (2002), 145--205.

\bibitem[Hochschild1965]{Hochschild1965}
G.~Hochschild,
\textit{The Structure of Lie Groups}.
Holden-Day, San Francisco, 1965.

\bibitem[Brown82]{Brown82}
K.~S. Brown,
\textit{Cohomology of Groups}.
Springer, 1982.

\end{thebibliography}

\end{document}



